\documentclass[11pt]{article}
\usepackage[a4paper,margin=1in]{geometry}
\usepackage{fontspec}
\usepackage[boldfont=false]{xeCJK}
\setCJKmainfont{FandolSong-Regular.otf}
\usepackage{amsmath}
\usepackage{amssymb}
\usepackage{array}
\usepackage{booktabs}
\usepackage{graphicx}
\usepackage{adjustbox}
\usepackage{enumitem}
\setlist[itemize]{topsep=2pt,partopsep=0pt,parsep=0pt,itemsep=2pt,leftmargin=1.4em}
\setlist[enumerate]{topsep=2pt,partopsep=0pt,parsep=0pt,itemsep=2pt,leftmargin=1.6em}
\usepackage{parskip}
\usepackage{fvextra}
\fvset{breaklines=true,breakanywhere=true,fontsize=\small}
\setlength{\parindent}{0pt}

\begin{document}

\begin{center}
  {\LARGE\bfseries Transport in plants}\\[0.35em]
  {\large A Level Biology}
\end{center}
\vspace{0.6em}

\section*{The two transport tissues}

Plants move substances through two transport \textbf{tissues} 组织:

\begin{itemize}
\item \textbf{xylem} 木质部 carries water and dissolved \textbf{mineral ions} 矿物离子 up from the roots.

\item \textbf{phloem} 韧皮部 carries dissolved foods (called \textbf{assimilates} 同化物, mainly sugars) to wherever they are needed.

\end{itemize}

In a \textbf{transverse section} 横切面 (a cut straight across) of a \textbf{dicotyledonous plant} 双子叶植物:

\begin{itemize}
\item in the \textbf{stem}, xylem and phloem sit together in bundles near the outside, with xylem on the inside of each bundle.

\item in the \textbf{root}, the xylem is in the centre, often in a star shape, with phloem between the arms.

\item in the \textbf{leaf}, both are found in the veins.

\end{itemize}

When you draw a plan diagram, you draw only the outlines of the tissues, not the single cells.

\subsection*{Xylem vessels}

Xylem water-carrying tubes are called \textbf{vessels} 导管. They are made of dead, empty cells joined end to end, with the end walls gone, so they form one long open pipe. Their walls are thickened and waterproofed with \textbf{lignin} 木质素. So the structure suits the job: the hollow, open tube with no contents lets water flow fast, and the lignin gives strength and support.

\subsection*{Phloem sieve tubes and companion cells}

Phloem food-carrying tubes are called \textbf{sieve tubes} 筛管. They are living cells joined end to end, but their end walls are not gone — they become \textbf{sieve plates} 筛板 with many holes that sap flows through. To leave room for flow, a sieve tube cell loses most of its contents and has no nucleus.

Beside each sieve tube is a \textbf{companion cell} 伴胞. It keeps its \textbf{nucleus} 细胞核 and has many \textbf{mitochondria} 线粒体. It does the living work for the sieve tube and loads sugars into it.

\section*{Water from the soil to the xylem}

Water enters a \textbf{root hair cell} 根毛细胞 by \textbf{osmosis} 渗透, because the root hair has a lower water potential than the soil water. Water then crosses the root to the xylem by two pathways:

\begin{itemize}
\item the \textbf{apoplast pathway} 质外体途径 — water moves through the \textbf{cell walls} 细胞壁 (made of \textbf{cellulose} 纤维素) and the spaces between cells, without entering the cytoplasm. This is fast.

\item the \textbf{symplast pathway} 共质体途径 — water moves through the cytoplasm of cells, passing from cell to cell through the \textbf{plasmodesmata} 胞间连丝.

\end{itemize}

At a ring of cells called the \textbf{endodermis} 内皮层, the apoplast pathway is blocked by the \textbf{Casparian strip} 凯氏带, a waterproof band of \textbf{suberin} 木栓质. This forces all the water through the cell membranes, which lets the plant control what enters the xylem.

\section*{Transpiration and the movement of water up the xylem}

\textbf{Transpiration} 蒸腾作用 is the loss of water vapour from a plant. Water \textbf{evaporates} (turns to vapour) from the wet cell surfaces inside the leaf — this is \textbf{evaporation} 蒸发. The \textbf{water vapour} 水蒸气 then \textbf{diffuses} 扩散 out through the \textbf{stomata} 气孔 into the \textbf{atmosphere} 大气.

This loss at the top pulls water up the xylem in a continuous column. It works because of hydrogen bonding between water molecules:

\begin{itemize}
\item water molecules attract each other through \textbf{hydrogen bonds} 氢键, so they stick together. This sticking is \textbf{cohesion} 内聚力, and it lets the whole column be pulled up under \textbf{tension} 张力 (the cohesion–tension idea).

\item water molecules also stick to the cellulose of the cell walls. This is \textbf{adhesion} 附着力, which helps hold the column in place.

\end{itemize}

\section*{Xerophytes}

A \textbf{xerophyte} 旱生植物 is a plant \textbf{adapted} 适应 to live where water is scarce. Its leaves reduce water loss by transpiration in several ways: a thick waxy \textbf{cuticle} 角质层, \textbf{stomata} sunk in pits, hairs that trap moist air, and leaves that can roll up. You should be able to draw a labelled leaf section showing these features.

\section*{Translocation: moving assimilates in the phloem}

Assimilates such as \textbf{sucrose} 蔗糖 and \textbf{amino acids} 氨基酸 are carried in the phloem from a \textbf{source} 源 to a \textbf{sink} 库.

\begin{itemize}
\item a \textbf{source} is where the assimilate is made or released (for example a photosynthesising leaf).

\item a \textbf{sink} is where it is used or stored (for example a growing root).

\end{itemize}

\subsection*{Loading at the source}

Companion cells load sucrose into the sieve tubes against its concentration gradient. They use \textbf{proton pumps} 质子泵 to pump hydrogen ions out, then \textbf{cotransporter proteins} 协同运输蛋白 bring sucrose back in together with those ions. This is a form of \textbf{active transport} 主动运输.

Loading sucrose lowers the \textbf{water potential} 水势 inside the sieve tube, so water follows by osmosis. This raises the \textbf{hydrostatic pressure} 静水压 there.

\subsection*{Mass flow}

At the sink, sucrose is removed, so the water potential rises, water leaves, and the pressure falls. The result is a pressure difference between source (high) and sink (low). Sap flows from high to low pressure down this gradient. This pressure-driven flow is called \textbf{mass flow} 集流.


\end{document}
